\naslov{Odločitvena drevesa}

Odločitveno drevo razbije prostor na dele, na katerih je napoved konstantna.
Pri tem uporabljamo le teste oblike $X_i < a_j$; ena veja predstavlja resnično,
ena pa neresnično vrednost tega izraza.
Končna vozlišča podajo napovedi.
Za konstrukcijo uporabimo algoritem TDIDT (\pojem{top-down induction of decision
  trees}).
Definiramo mero nečistoče (\pojem{impurity}) za neko množico; potem je nečistoča
drevesa enaka
\[
  \text{Impurity}(T) = \sum_{l \in T} \frac{\abs{S_l}}{\abs{S}}
  \text{Impurity}(S_l),
\]
kjer so $l$ listi drevesa.
Želimo poiskati drevo z najmanjšo nečistočo, to pa je žal NP-poln problem, zato
uporabimo požrešni algoritem; v vsakem koraku vzamemo poddrevesi, ki imata
skupaj najmanjšo nečistočo.
Taki poddrevesi poiščemo s polnim iskanjem možnih testov.
Ko napovedujemo vrednosti, je napoved lista enaka povprečju vrednosti ciljne
spremenljivke v tem listu (ali večinska vrednost v diskretnem primeru).

\vprasanje{Kako deluje algoritem TDIDT?}

Pri izbiri mere nečistoče ločimo numerične in diskretne primere.
V numeričnem merimo varianco
\[
  \text{Impurity}(S) = \inv{\abs{S}} \sum_{(\mb{x}, y) \in S} \left( y - \bar{y}
  \right)^2,
\]
pri diskretnem primeru pa je $\text{Impurity}(S) = \phi(p_1, \ldots, p_c)$,
kjer so
\[
  p_i = P(Y = y_i \such S) = \frac{\abs{\{y = y_i\}}}{\abs{S}},
\]
$\phi$ pa je neka funkcija, za katero želimo, da čim bolje modelira varianco.
Zahtevamo, da ima največjo vrednost pri enakomerni porazdelitvi, minimalno pri
determinističnih, ter da je neodvisna od permutacije argumentov.
Imamo dve pogosti možnosti:
\begin{itemize}
\item Entropija: $\phi = \sum_i p_i \log_2 p_i$
\item Indeks Gini: $\phi = 1 - \sum_i p_i^2$
\end{itemize}

\vprasanje{Kakšne mere nečistoče poznamo pri gradnji odločitvenih dreves?}

Če so drevesa prevelika, lahko v listih dovoliš $\text{Impurity} > 0$, čemur
pravimo \pojem{tree pruning}.
To lahko naredimo na dva načina.
Če si nastavimo minimalno število primerov, ki jih ne bomo več ločevali, temu
pravimo \pojem{sprotno rezanje}, če pa prvo naredimo celo drevo, in nato neka
notranja vozlišča spremenimo v končna, pa govorimo o \pojem{naknadnem rezanju}.
Pri tem gledamo napako poddrevesa
\[
  \text{Err}_\alpha(T) = \text{Err}(T) + \alpha \abs{T},
\]
kjer je $T$ število listov v drevesu, $\alpha$ pa izbran parameter.
V algoritmu primerjamo napako v staršu in v otrocih, in vzamemo manjšo.

\vprasanje{Kako upravljaš velikost odločitvenega drevesa?}